\chapter{実験評価}
\label{chap:evaluation}

\section{評価目的と実験条件}
\subsection{評価項目と比較対象}
\writehere{検証したい仮説，各評価指標，比較手法，実験回数を定義する．}
\subsection{実験環境と機材}
\writehere{実施場所，期間，機材，ソフトウェア版，モデル，閾値，電源条件を記録する．}
\subsection{評価データと正解ラベル}
\writehere{データの収集・分割方法，昼夜の区分，対象数，正解ラベルの作成手順を記す．画像数と撮影サイクル数を区別する．}

\section{動物検出性能}
\writehere{評価単位と混同行列を定義し，適合率，再現率，見逃しを昼夜別・条件別に評価する．モデル出力を参照ラベルとする場合は，人手正解との違いを明示する．}

\section{画像選別と通信量削減}
\writehere{画像数，撮影サイクル数，転送バイト数について，削減率の分母と分子を定義して結果を示す．}

\section{通信性能と通知遅延}
\writehere{距離，障害物，電波強度，転送成功率，撮影から通知までの時間を測定する．測定区間と時刻同期方法を示す．}

\section{消費電力と連続稼働時間}
\writehere{測定点，電圧，各動作状態の電流と時間を記録し，実測結果と電池寿命の試算を区別する．}

\section{導入・運用コスト}
\writehere{機材費，回線費，電源設備，保守作業を同一条件で比較する．価格の調査時点を明記する．}

\section{評価結果のまとめ}
\writehere{各結果を研究目的と対応付け，確認できた事項と未検証の事項を整理する．}
